\section{Three-modulus face calculus and determinant-first reassembly}
\label{sec:tpc37-three-modulus}

After the physical zero-coordinate excision of
\cref{sec:tpc37-physical-zero-faces}, the target \(N\) and both factors
in \(su=N\) are units modulo each auxiliary prime.  This section keeps
three orbit-scale primes separate and computes all eight multiplicative
character faces exactly.  The calculation has two consequences.  First,
the seven proper-coordinate faces have one extra power of \(J^{-1}\) in
quadratic mass.  Second, grouping the completed Cartesian block by the
integer determinant \(F=su-N\), rather than completing \(s\) or \(u\)
separately, removes the apparent incomplete-box boundary loss.  Neither
statement is a bound for the physical coherent row energy; the precise
remaining Bessel-type input is isolated at the end.

\subsection{The eight faces on the correct finite-field domain}

Let \(q_1,q_2,q_3>2\) be distinct primes, put
\begin{equation}
 n_i=q_i-1,
 \qquad p_i=\frac1{n_i},
 \qquad M=q_1q_2q_3,
 \label{eq:tpc37-three-modulus-data}
\end{equation}
and fix an integer \(N\) satisfying
\begin{equation}
 (N,M)=1.
 \label{eq:tpc37-N-unit}
\end{equation}
For \(F\in\mathbb Z/q_i\mathbb Z\), define
\begin{align}
 d_i(F)&:=\mathbf 1_{F=0},
 \label{eq:tpc37-coordinate-delta}\\
 e_i(F)&:=\mathbf 1_{F\ne-N},
 \label{eq:tpc37-coordinate-unit-mask}\\
 m_i(F)&:=p_i e_i(F),
 \qquad
 c_i(F):=d_i(F)-m_i(F).
 \label{eq:tpc37-principal-centered-coordinate}
\end{align}
The exclusion \(F=-N\) in \(e_i\) is essential.  Indeed, if
\(F=su-N\), then \(F=-N\pmod {q_i}\) is exactly the residue on which
\(su=0\pmod {q_i}\), so every multiplicative character extended by zero
must vanish there.

To make the character interpretation explicit, write
\begin{equation}
 x_i(F):=(N+F)N^{-1}\pmod {q_i}.
 \label{eq:tpc37-ratio-coordinate}
\end{equation}
Every multiplicative character of
\(G_i:=\mathbb F_{q_i}^{\times}\) is extended by zero at the origin.  The
orthogonality relation on \(G_i\) gives
\begin{align}
 d_i(F)
 &=\frac1{n_i}\sum_{\chi_i}\chi_i(x_i(F)),
 \label{eq:tpc37-all-character-coordinate}\\
 m_i(F)
 &=\frac1{n_i}\mathbf 1_{G_i}(x_i(F)),
 \label{eq:tpc37-principal-character-coordinate}\\
 c_i(F)
 &=\frac1{n_i}\sum_{\chi_i\ne\mathbf 1}
 \chi_i(x_i(F)).
 \label{eq:tpc37-nonprincipal-character-coordinate}
\end{align}
Thus \(m_i\) is the principal multiplicative face and \(c_i\) is the
aggregate nonprincipal face.  In particular,
\begin{equation}
 d_i=m_i+c_i
 \label{eq:tpc37-one-coordinate-face-split}
\end{equation}
holds also at \(x_i=0\), where all three terms vanish.

For \(S\subseteq[3]:=\{1,2,3\}\), define the corresponding face on
\(\mathbb Z/M\mathbb Z\) by
\begin{equation}
 K_S(F)
 :=\prod_{i\in S}c_i(F)
   \prod_{i\notin S}m_i(F).
 \label{eq:tpc37-character-face}
\end{equation}
The CRT identifies \(F\pmod M\) with its three coordinates.  Therefore
\eqref{eq:tpc37-one-coordinate-face-split} gives the exact eight-face
identity
\begin{equation}
 \boxed{
 \mathbf 1_{M\mid F}
 =\prod_{i=1}^3d_i(F)
 =\sum_{S\subseteq[3]}K_S(F).}
 \label{eq:tpc37-eight-face-identity}
\end{equation}
All norms below are unnormalized counting norms on
\(\mathbb Z/M\mathbb Z\).

\begin{theorem}[Orthogonal face calculus]
\label{thm:tpc37-orthogonal-face-calculus}
The eight functions \(K_S\), \(S\subseteq[3]\), are pairwise orthogonal
in \(\ell^2(\mathbb Z/M\mathbb Z)\).  For every \(S\subseteq[3]\),
\begin{align}
 \boxed{
 \|K_S\|_2^2
 =\prod_{i\in S}(1-p_i)
  \prod_{i\notin S}p_i,}
 \label{eq:tpc37-face-l2}\\
 \boxed{
 \|K_S\|_1
 =2^{|S|}\prod_{i\in S}(1-p_i).}
 \label{eq:tpc37-face-l1}
\end{align}
Moreover, at the joint residue \(F=0\pmod M\),
\begin{equation}
 \boxed{
 K_S(0)=\|K_S\|_2^2.}
 \label{eq:tpc37-face-joint-value}
\end{equation}
Consequently the squared norms of the empty, singleton, double, and full
faces have respective scales
\begin{equation}
 J^{-3},\qquad J^{-2},\qquad J^{-1},\qquad 1
 \label{eq:tpc37-face-codimension-scales}
\end{equation}
when \(q_i\asymp J\).
\end{theorem}

\begin{proof}
In one coordinate,
\begin{equation}
 \sum_{F\bmod q_i}c_i(F)=0,
 \qquad
 \langle c_i,m_i\rangle=0.
 \label{eq:tpc37-one-coordinate-orthogonality}
\end{equation}
There are \(n_i\) residues on which \(e_i=1\).  Hence
\begin{align}
 \|m_i\|_2^2&=n_i p_i^2=p_i,
 &\|m_i\|_1&=n_i p_i=1,
 \label{eq:tpc37-principal-coordinate-norms}\\
 \|c_i\|_2^2
 &=(1-p_i)^2+(n_i-1)p_i^2=1-p_i,
 &\|c_i\|_1
 &=(1-p_i)+(n_i-1)p_i=2(1-p_i).
 \label{eq:tpc37-centered-coordinate-norms}
\end{align}
Tensoring through the CRT proves orthogonality and
\eqref{eq:tpc37-face-l2}--\eqref{eq:tpc37-face-l1}.  At \(F=0\), every
coordinate equals the distinguished unit residue, so
\begin{equation}
 c_i(0)=1-p_i,
 \qquad m_i(0)=p_i,
 \label{eq:tpc37-coordinate-joint-values}
\end{equation}
which proves \eqref{eq:tpc37-face-joint-value}.  The scale table follows
from \(p_i\asymp J^{-1}\).
\end{proof}

\subsection{The coherent union of all proper faces}

Put
\begin{equation}
 A:=\prod_{i=1}^3(1-p_i),
 \qquad
 C(F):=K_{[3]}(F)=\prod_{i=1}^3c_i(F),
 \label{eq:tpc37-full-face-A}
\end{equation}
and combine the seven proper faces before taking absolute values:
\begin{equation}
 P(F):=\sum_{S\subsetneq[3]}K_S(F).
 \label{eq:tpc37-proper-union-definition}
\end{equation}
By \eqref{eq:tpc37-eight-face-identity},
\begin{equation}
 \boxed{P(F)=\mathbf 1_{M\mid F}-C(F).}
 \label{eq:tpc37-proper-union-identity}
\end{equation}
In particular, on every nonjoint residue the proper union and full face
cancel exactly.

\begin{theorem}[Exact norms of the proper union]
\label{thm:tpc37-proper-union-norms}
With the preceding definitions,
\begin{align}
 \boxed{\|P\|_\infty=1-A,}
 \label{eq:tpc37-proper-linf}\\
 \boxed{\|P\|_2^2=1-A,}
 \label{eq:tpc37-proper-l2}\\
 \boxed{\|P\|_1=1+6A.}
 \label{eq:tpc37-proper-l1}
\end{align}
The full face satisfies
\begin{equation}
 \|C\|_2^2=A,
 \qquad
 \|C\|_1=8A.
 \label{eq:tpc37-full-face-norms}
\end{equation}
Thus, for \(q_i\asymp J\),
\begin{equation}
 1-A\asymp J^{-1}.
 \label{eq:tpc37-proper-one-over-J}
\end{equation}
\end{theorem}

\begin{proof}
The full-face formulas are the \(S=[3]\) case of
\cref{thm:tpc37-orthogonal-face-calculus}.  Orthogonality of all eight
faces and \eqref{eq:tpc37-eight-face-identity} give
\begin{equation}
 \|P\|_2^2
 =\sum_{S\subsetneq[3]}\|K_S\|_2^2
 =1-A.
 \label{eq:tpc37-proper-l2-proof}
\end{equation}
At the joint residue, \(P(0)=1-A\).  Away from the joint residue, either
one coordinate is the excluded zero residue, in which case \(P=C=0\),
or \eqref{eq:tpc37-proper-union-identity} gives \(P=-C\).  If \(Z\) is
the nonempty set of failed congruence coordinates at such a residue,
then
\begin{equation}
 |C(F)|
 =\prod_{i\in Z}p_i
  \prod_{i\notin Z}(1-p_i)
 \le \max_i p_i
 \le1-A.
 \label{eq:tpc37-proper-pointwise-failed}
\end{equation}
This proves \eqref{eq:tpc37-proper-linf}.

Finally, the joint residue contributes \(1-A\) to \(\|P\|_1\), while
the nonjoint residues contribute
\begin{equation}
 \|C\|_1-|C(0)|=8A-A=7A.
 \label{eq:tpc37-proper-offjoint-l1}
\end{equation}
Their sum is \(1+6A\), proving \eqref{eq:tpc37-proper-l1}.  The final
comparison follows from fixedly many \(p_i\asymp J^{-1}\).
\end{proof}

The same identities locate the quadratic mass exactly.

\begin{corollary}[Joint-alias and ghost-face energy]
\label{cor:tpc37-joint-ghost-energy}
On one period modulo \(M\),
\begin{align}
 |C(0)|^2&=A^2,
 &|P(0)|^2&=(1-A)^2,
 \label{eq:tpc37-joint-face-energy}\\
 \sum_{\substack{F\bmod M\\F\ne0}}|C(F)|^2
 &=A(1-A),
 &\sum_{\substack{F\bmod M\\F\ne0}}|P(F)|^2
 &=A(1-A).
 \label{eq:tpc37-ghost-face-energy}
\end{align}
Thus the full-face energy is concentrated on joint aliases, whereas the
proper-face energy is concentrated on nonjoint residues that cancel the
full face in \eqref{eq:tpc37-proper-union-identity}.
\end{corollary}

\begin{proof}
The joint values follow from \eqref{eq:tpc37-full-face-A} and
\eqref{eq:tpc37-proper-union-identity}.  Subtract them from the total
squared norms in \eqref{eq:tpc37-proper-l2} and
\eqref{eq:tpc37-full-face-norms}.
\end{proof}

\subsection{Determinant-first periodization}

The finite calculation becomes useful on the physical reciprocal blocks
only after choosing the correct integer variable.  Fix \(0<N\ll X\)
satisfying \eqref{eq:tpc37-N-unit}, and let
\(\mathcal R\subset\mathbb N^2\) be any set of pairs \((s,u)\) such that
\begin{equation}
 |su-N|\le B,
 \qquad B\ll X.
 \label{eq:tpc37-determinant-range}
\end{equation}
TPC-36 supplies this hypothesis on every reciprocal dyadic block.  Put
\begin{equation}
 F(s,u)=su-N.
 \label{eq:tpc37-integer-determinant}
\end{equation}

\begin{theorem}[Determinant-first periodization]
\label{thm:tpc37-determinant-periodization}
Let \(H\) be any of the eight faces \(K_S\), the full face \(C\), or the
proper union \(P\).  For \(r\in\{1,2\}\) and every fixed
\(\varepsilon>0\),
\begin{equation}
 \boxed{
 \sum_{(s,u)\in\mathcal R}|H(F(s,u))|^r
 \ll_{\varepsilon}
 X^{\varepsilon}
 \left(1+\frac BM\right)
 \sum_{F\bmod M}|H(F)|^r.}
 \label{eq:tpc37-determinant-periodization-bound}
\end{equation}
The same estimate holds after inserting unary weights of pointwise size
\(X^{o(1)}\).  It is uniform in the shape of the reciprocal rectangle.
\end{theorem}

\begin{proof}
For a fixed integer \(F\), every pair counted on the left satisfies
\begin{equation}
 su=N+F.
 \label{eq:tpc37-fixed-determinant-factorization}
\end{equation}
When \(N+F>0\), the number of positive representations is at most
\(\tau(N+F)\ll_\varepsilon X^\varepsilon\); when \(N+F\le0\), there are
none.  Therefore
\begin{equation}
 \sum_{(s,u)\in\mathcal R}|H(F(s,u))|^r
 \ll_\varepsilon X^\varepsilon
 \sum_{|F|\le B}|H(F)|^r.
 \label{eq:tpc37-divisor-multiplicity-reduction}
\end{equation}
The function \(H\) is periodic modulo \(M\).  Partitioning
\([-B,B]\cap\mathbb Z\) into complete residue blocks and at most two
partial blocks gives
\begin{equation}
 \sum_{|F|\le B}|H(F)|^r
 \ll\left(1+\frac BM\right)
 \sum_{F\bmod M}|H(F)|^r.
 \label{eq:tpc37-period-block-bound}
\end{equation}
Combining the last two displays proves the theorem.  Pointwise-soft unary
weights are absorbed into \(X^\varepsilon\).
\end{proof}

The theorem is the reason to group by \(F\) before completing either
\(s\) or \(u\).  A one-variable residue count produces an artificial
boundary such as \(\min\{S,U\}/J\); the determinant fiber has only
divisor-bounded multiplicity and sees the exact \(M\)-periodic norm.

\begin{corollary}[Endpoint periodization ledger]
\label{cor:tpc37-periodization-ledger}
Assume \(q_i\asymp J\), \(B\ll X\), and put
\begin{equation}
 \mathcal K_X:=1+\frac X{J^3}.
 \label{eq:tpc37-critical-period-count}
\end{equation}
Then every individual face and the proper union satisfy
\begin{equation}
 \sum_{(s,u)\in\mathcal R}|K_S(F(s,u))|
 +\sum_{(s,u)\in\mathcal R}|P(F(s,u))|
 \ll X^{o(1)}\mathcal K_X.
 \label{eq:tpc37-all-face-l1-periodization}
\end{equation}
In quadratic mass,
\begin{align}
 \sum_{(s,u)\in\mathcal R}|P(F(s,u))|^2
 &\ll X^{o(1)}\frac{\mathcal K_X}{J},
 \label{eq:tpc37-proper-l2-periodization}\\
 \sum_{(s,u)\in\mathcal R}|C(F(s,u))|^2
 &\ll X^{o(1)}\mathcal K_X.
 \label{eq:tpc37-full-l2-periodization}
\end{align}
At the endpoint \eqref{eq:tpc37-endpoint-scales},
\begin{align}
 \mathcal K_X
 &=\frac X{J^3}X^{o(1)}
 =\frac Q{J^2}X^{o(1)}
 =X^{1/400+o(1)},
 \label{eq:tpc37-critical-period-endpoint}\\
 \frac{\mathcal K_X}{J}
 &=\frac X{J^4}X^{o(1)}
 =X^{-132/400+o(1)}.
 \label{eq:tpc37-proper-period-decay}
\end{align}
\end{corollary}

\begin{proof}
Apply \cref{thm:tpc37-determinant-periodization} with the exact norms in
\cref{thm:tpc37-orthogonal-face-calculus,thm:tpc37-proper-union-norms}.
Since every face has \(O(1)\) one-period \(\ell^1\) norm, the first bound
follows.  The proper and full \(\ell^2\) bounds use respectively
\(1-A\asymp J^{-1}\) and \(A\asymp1\).  The endpoint identities follow
from \(QJ=X^{1+o(1)}\) and \(J=X^{133/400+o(1)}\).
\end{proof}

\begin{remark}[Quadratic sparsity is not physical closure]
\label{rem:tpc37-quadratic-not-physical}
Equation \eqref{eq:tpc37-proper-period-decay} is an entrywise/fiberwise
quadratic statement.  The physical orbit slice sums in \(s,u\), the input
rows, and the determinant shifts before taking its square.  No inequality
in this section transfers \eqref{eq:tpc37-proper-period-decay} through
those coherent sums.  In particular, it is not a proof of the inherited
coefficient gate.
\end{remark}

\subsection{Aliases, proper-coordinate ghosts, and the terminal fiber}

For every integer \(k\), the joint alias \(F=kM\) has the same face values
as \(F=0\):
\begin{equation}
 C(kM)=A,
 \qquad
 P(kM)=1-A.
 \label{eq:tpc37-joint-alias-values}
\end{equation}
Thus the \(O(X/M)\) joint aliases have respective quadratic masses
\begin{equation}
 O\!\left(\frac XM\right)
 \quad\text{and}\quad
 O\!\left(\frac X{MJ^2}\right)
 \label{eq:tpc37-joint-alias-mass-ledger}
\end{equation}
on the full and proper faces.  The larger proper total in
\eqref{eq:tpc37-proper-l2-periodization} comes from nonjoint residues
where \(P=-C\).

This cancellation is not optional.  For example, work in a terminal-size
interval of \(u\)'s of length \(\asymp X\), set \(s=1\), and suppose
\((N,M)=1\).  On the family
\begin{equation}
 u=N+kq_1q_2,
 \qquad
 q_3\nmid k,
 \qquad
 q_3\nmid N+kq_1q_2,
 \label{eq:tpc37-two-hit-ghost-family}
\end{equation}
the first two congruences hold, the third fails, and all three ratio
coordinates are nonzero.  There are
\begin{equation}
 \asymp\frac X{q_1q_2}\asymp\frac X{J^2}
 \label{eq:tpc37-two-hit-count}
\end{equation}
such points.  At each one,
\begin{equation}
 C(F)=-(1-p_1)(1-p_2)p_3,
 \qquad P(F)=-C(F),
 \label{eq:tpc37-two-hit-cancellation}
\end{equation}
so the full face alone has total absolute mass
\begin{equation}
 \asymp\frac X{J^2}\frac1J
 =\frac X{J^3}.
 \label{eq:tpc37-two-hit-critical-mass}
\end{equation}
The original joint CRT delta is zero on the whole family.  Dropping the
proper faces before their cancellation with \(C\) would therefore create
a false alias family at exactly the critical three-modulus scale.

On the genuine integer equality \(F=0\), the situation is different.
Every proper divisor layer and the terminal layer \(s=1,u=N\) acquire the
same scalar face values
\begin{equation}
 P(0)=1-A\ll J^{-1},
 \qquad C(0)=A=1-O(J^{-1}).
 \label{eq:tpc37-exact-equality-face-values}
\end{equation}
Let \(V_{\rm eq}^{\rm abs}\) denote the sign-blind physical energy of the
regular exact-equality kernel.  The inherited envelopes give
\begin{equation}
 V_{\rm eq}^{\rm abs}\ll X^{o(1)}Q^3J.
 \label{eq:tpc37-equality-absolute-baseline}
\end{equation}
The exact scalar relation in
\eqref{eq:tpc37-exact-equality-face-values} therefore yields
\begin{equation}
 \boxed{
 V_{{\rm eq},{\rm proper}}
 \ll X^{o(1)}J^{-2}Q^3J
 =X^{o(1)}\frac{Q^3}{J}.}
 \label{eq:tpc37-equality-proper-closure}
\end{equation}
This closes all proper faces on the actual equality without using a sign
estimate.  By contrast, the full face is \(A\) times the same vector and
retains it with asymptotic coefficient one.  In particular, the fused
terminal M\"obius correlation is not weakened by coordinate centering.

\begin{corollary}[Full-face centering obstruction]
\label{cor:tpc37-full-face-obstruction}
No uniform power saving for vectors supported on \(F=0\) can follow from
three-coordinate multiplicative centering alone.  More precisely, on
that subspace the full-face operator is scalar multiplication by
\begin{equation}
 A=1-O(J^{-1}).
 \label{eq:tpc37-full-face-scalar}
\end{equation}
This is a method obstruction for the CRT-centering route, not a lower
bound for the physical M\"obius correlation.
\end{corollary}

\subsection{A sharp coherent no-go and the next gate}

The \(L^2\) saving for \(P\) cannot be upgraded uniformly over all
bounded unary coefficients.  Retain \(q_i\asymp J\) and the endpoint
scales.  In the same terminal geometry, take \(s=1\), let \(I\) be an
interval of \(u\)'s of length \(B\asymp X\), and assume \(B/M\to\infty\).
Define the bounded unary weight
\begin{equation}
 g_N(u)=
 \begin{cases}
  \overline{P(u-N)}/|P(u-N)|,&P(u-N)\ne0,\\
  0,&P(u-N)=0.
 \end{cases}
 \label{eq:tpc37-coherent-saturating-weight}
\end{equation}

\begin{proposition}[Sharp quadratic-to-coherent separation]
\label{prop:tpc37-coherent-face-no-go}
For the preceding weight,
\begin{align}
 \left|\sum_{u\in I}g_N(u)P(u-N)\right|
 &\asymp\frac BM,
 \label{eq:tpc37-coherent-saturation-l1}\\
 \sum_{u\in I}|P(u-N)|^2
 &\asymp\frac B{MJ}.
 \label{eq:tpc37-coherent-saturation-l2}
\end{align}
At the endpoint \(B\asymp X\), the ratio of the first quantity to the
square root of the second is
\begin{equation}
 \boxed{
 \asymp\sqrt{\mathcal K_X J}
 =\sqrt{\frac X{J^2}}
 =X^{67/400+o(1)}.}
 \label{eq:tpc37-coherent-upgrade-loss}
\end{equation}
Thus no estimate uniform over all bounded unary coefficients can upgrade
the displayed periodic \(L^2\) mass to the required coherent scale with
only an \(X^{o(1)}\) loss.
\end{proposition}

\begin{proof}
The definition of \(g_N\) makes the first left side equal to
\(\sum_{u\in I}|P(u-N)|\).  Decompose \(I-N\) into complete \(M\)-periods
and at most two partial periods.  The exact one-period norms
\eqref{eq:tpc37-proper-l1}--\eqref{eq:tpc37-proper-l2}, together with
\(B/M\to\infty\), give
\eqref{eq:tpc37-coherent-saturation-l1}--
\eqref{eq:tpc37-coherent-saturation-l2}.  Finally use
\(\mathcal K_X=X/J^3\) and the endpoint value of \(J\).
\end{proof}

The saturating \(g_N\) is an artificial periodic coefficient.  The
proposition does not assert that the physical weight
\(-\mu(u)\log u\) has this alignment.  It proves that avoiding it requires
coefficient-specific arithmetic information.

To state the remaining target correctly, the dependence on the moving
physical target must be retained.  Restore it in the notation by writing
\(C_N(F)\) and \(P_N(F)\).  The one-period norms are independent of the
unit \(N\bmod M\).  Therefore, directly from periodicity and the exact
norms---without using divisor multiplicity---one has, uniformly for
\((N,M)=1\),
\begin{align}
 \sum_{|F|\ll X}|C_N(F)|^2
 &\ll\mathcal K_X,
 \notag\\
 \sum_{|F|\ll X}|P_N(F)|^2
 &\ll\frac{\mathcal K_X}{J}.
 \label{eq:tpc37-direct-face-period-ledger}
\end{align}

Let \(\mathcal N_{\rm reg}\) denote the targets occurring in one physical
packet after the zero-face excision, so every
\(N\in\mathcal N_{\rm reg}\) satisfies \((N,M)=1\).  Let
\begin{equation}
 \mathscr H_{\rm phys}
 :=\ell^2\{(L,\gamma,j)\}
   \oplus\ell^2\{(R,\alpha,j)\}
 \label{eq:tpc37-physical-output-Hilbert-space}
\end{equation}
be the direct sum of the left and right packet output spaces, with the
indices restricted to the retained packet.  For the unweighted physical
output its squared norm is exactly \(V_L+V_R\).  All dyadic, projective,
and reassembly layers belonging to that packet, including their proved
\(X^{o(1)}\) cost, are included in the outputs below.
For a scalar field
\(\boldsymbol\kappa=(\kappa_N(F))_{N\in\mathcal N_{\rm reg}}\), let
\(\mathscr Z_F[\boldsymbol\kappa]\) denote the Hilbert-space-valued
physical output of the shifted determinant fiber
\begin{equation}
 su=\ell dj+h+F.
 \label{eq:tpc37-shifted-determinant-fiber}
\end{equation}
Thus \(\mathscr Z_F[\boldsymbol\kappa]\in\mathscr H_{\rm phys}\), and
every atomic summand with target \(N=\ell dj+h\) is multiplied by
\(\kappa_N(F)\) before the coherent row sum.  This notation is necessary:
the coordinate faces depend on \(N\bmod M\) and cannot be pulled out as a
single scalar depending only on \(F\).

\begin{assumption}[Target-coupled shifted-fiber Bessel input]
\label{ass:tpc37-shifted-fiber-bessel}
For the actual row mask and M\"obius-derived coefficients,
\begin{equation}
 \boxed{
 \left\|\sum_{|F|\ll X}
 \mathscr Z_F[\boldsymbol\kappa]\right\|_2^2
 \ll X^{o(1)}Q^2J
 \sup_{N\in\mathcal N_{\rm reg}}
 \sum_{|F|\ll X}|\kappa_N(F)|^2}
 \label{eq:tpc37-shift-fiber-bessel-gate}
\end{equation}
for the two fields \(\kappa_N=C_N\) and \(\kappa_N=P_N\).
\end{assumption}

This is a sufficient input, not a theorem of the present paper.  Under
\cref{ass:tpc37-shifted-fiber-bessel},
\eqref{eq:tpc37-direct-face-period-ledger} gives
\begin{equation}
 Q^2J\mathcal K_X
 =Q^2J\frac Q{J^2}
 =\frac{Q^3}{J},
 \label{eq:tpc37-full-face-bessel-ledger}
\end{equation}
for the full face, and the stronger scale
\begin{equation}
 Q^2J\frac{\mathcal K_X}{J}
 =\frac{Q^3}{J^2}.
 \label{eq:tpc37-proper-face-bessel-ledger}
\end{equation}

The conditional input
\eqref{eq:tpc37-shift-fiber-bessel-gate} is not a formal consequence of
the periodic \(L^2\) calculation.
\Cref{prop:tpc37-coherent-face-no-go} shows that arbitrary bounded unary
data can align with the \(L^1\), rather than \(L^2\), size of the face.
Any proof at the stated physical normalization must therefore exploit
the actual row and M\"obius signs, or an equivalent dispersion
mechanism.  There is also a genuine uniformity issue: on the \(F\)-fiber
the affine shift is
\begin{equation}
 h_F=h+F,
 \label{eq:tpc37-moving-shift}
\end{equation}
which may be of size \(X\).  Fixed-\(h\) primitivity and incidence bounds
from the preceding papers cannot simply be invoked uniformly in \(F\).
The three-modulus calculation therefore removes the fixed-target
proper-coordinate and incomplete-period bookkeeping obstructions, but
leaves a precise target-coupled, coefficient-weighted shifted-fiber gate.
